\documentclass[12pt]{article}
\usepackage[a4paper,margin=1in]{geometry}
\usepackage{amsmath, amssymb}
\usepackage{graphicx}
\usepackage{xcolor}
\usepackage{tcolorbox}
\usepackage{multicol}
\usepackage{titlesec}
\usepackage{enumitem}
\setlength{\columnsep}{1cm}

% Section title styles
\titleformat{\section}{\Large\bfseries\sffamily\color{blue!70!black}}{}{0em}{}
\titleformat{\subsection}{\large\bfseries\color{purple!70!black}}{}{0em}{}

% tcolorbox styling
\tcbset{
  colback=blue!5!white,
  colframe=blue!75!black,
  boxrule=0.5pt,
  arc=3pt,
  left=4pt,
  right=4pt,
  top=4pt,
  bottom=4pt,
  fonttitle=\bfseries,
  title filled=true,
  coltitle=black,
}

\begin{document}


\begin{multicols}{2}

\section*{1. Arithmetic Basics}

\begin{tcolorbox}[title=Percent and Decimal Conversion]
\begin{itemize}[left=0pt]
  \item Percent to Decimal: \( x\% = \frac{x}{100} \)
  \item Decimal to Percent: \( \text{decimal} \times 100\% \)
  \item Increase/Decrease: \( \text{Original} \times (1 \pm \text{percent}) \)
\end{itemize}
\end{tcolorbox}

\begin{tcolorbox}[title=Ratios and Proportions]
\[
\frac{a}{b} = \frac{c}{d} \Rightarrow ad = bc
\]
\end{tcolorbox}

\begin{tcolorbox}[title=Average Formula]
\[
\text{Average} = \frac{\text{Sum of values}}{\text{Number of values}}
\]
\end{tcolorbox}

\section*{2. Algebra Essentials}

\begin{tcolorbox}[title=Algebraic Identities]
\[
(a + b)^2 = a^2 + 2ab + b^2
\]
\[
(a - b)^2 = a^2 - 2ab + b^2
\]
\[
(a + b)(a - b) = a^2 - b^2
\]
\end{tcolorbox}

\begin{tcolorbox}[title=Solving Linear Equations]
\[
ax + b = c \Rightarrow x = \frac{c - b}{a}
\]
\end{tcolorbox}

\section*{3. Geometry: Volume \& Surface Area}

\begin{tcolorbox}[title=Cube]
\[
V = s^3, \quad SA = 6s^2
\]
\end{tcolorbox}

\begin{tcolorbox}[title=Rectangular Prism]
\[
V = lwh, \quad SA = 2(lw + lh + wh)
\]
\end{tcolorbox}

\begin{tcolorbox}[title=Sphere]
\[
V = \frac{4}{3}\pi r^3, \quad SA = 4\pi r^2
\]
\end{tcolorbox}

\begin{tcolorbox}[title=Cylinder]
\[
V = \pi r^2 h, \quad SA = 2\pi r^2 + 2\pi rh
\]
\end{tcolorbox}

\begin{tcolorbox}[title=Cone]
\[
V = \frac{1}{3}\pi r^2 h, \quad SA = \pi r^2 + \pi r\ell
\]
\textit{Note: } \( \ell \) is the slant height
\end{tcolorbox}

\begin{tcolorbox}[title=Square Pyramid]
\[
V = \frac{1}{3}b^2 h, \quad SA = b^2 + 2b\ell
\]
\end{tcolorbox}



\section*{4. Probability and Sequences}

\begin{tcolorbox}[title=Probability]
\[
P(\text{event}) = \frac{\text{favorable outcomes}}{\text{total outcomes}}
\]
\end{tcolorbox}

\begin{tcolorbox}[title=Permutations and Combinations]
Permutations: \( P(n, r) = \frac{n!}{(n - r)!} \)\\
Combinations: \( C(n, r) = \frac{n!}{r!(n - r)!} \)
\end{tcolorbox}

\begin{tcolorbox}[title=Arithmetic Sequence]
\[
a_n = a_1 + (n - 1)d
\]
\end{tcolorbox}

\section*{5. Unit Conversions}

\begin{tcolorbox}[title=Length (U.S. Customary)]
\begin{itemize}[left=0pt]
  \item \( 1 \text{ foot (ft)} = 12 \text{ inches (in)} \)
  \item \( 1 \text{ yard (yd)} = 3 \text{ ft} = 36 \text{ in} \)
  \item \( 1 \text{ mile (mi)} = 5,280 \text{ ft} = 1,760 \text{ yd} \)
\end{itemize}
\end{tcolorbox}

\begin{tcolorbox}[title=Weight (U.S. Customary)]
\begin{itemize}[left=0pt]
  \item \( 1 \text{ pound (lb)} = 16 \text{ ounces (oz)} \)
  \item \( 1 \text{ ton} = 2,000 \text{ lb} \)
\end{itemize}
\end{tcolorbox}

\begin{tcolorbox}[title=Liquid Volume (U.S. Customary)]
\begin{itemize}[left=0pt]
  \item \( 1 \text{ cup} = 8 \text{ fl oz} \)
  \item \( 1 \text{ pint (pt)} = 2 \text{ cups} \)
  \item \( 1 \text{ quart (qt)} = 2 \text{ pt} = 4 \text{ cups} \)
  \item \( 1 \text{ gallon (gal)} = 4 \text{ qt} = 8 \text{ pt} = 16 \text{ cups} \)
\end{itemize}
\end{tcolorbox}

\begin{tcolorbox}[title=Time]
\begin{itemize}[left=0pt]
  \item \( 1 \text{ minute} = 60 \text{ seconds} \)
  \item \( 1 \text{ hour} = 60 \text{ minutes} = 3,600 \text{ seconds} \)
  \item \( 1 \text{ day} = 24 \text{ hours} \)
  \item \( 1 \text{week} = 7 \text{ days} \)
  \item \( 1 \text{year} = 12 \text{ months} = 365 \text{ days} \)
\end{itemize}
\end{tcolorbox}

\begin{tcolorbox}[title=Metric System]
\begin{itemize}[left=0pt]
  \item Length: 
    \[
    1 \text{ kilometer (km)} = 1,000 \text{ meters (m)}
    \]
    \[
    1 \text{ meter} = 100 \text{ (cm)} 
    \]
  \item Weight:
    \[
    1 \text{ kilogram (kg)} = 1,000 \text{ grams (g)}
    \]
  \item Volume:
    \[
    1 \text{ liter (L)} = 1,000 \text{ milliliters (mL)}
    \]
\end{itemize}
\end{tcolorbox}

\begin{tcolorbox}[title=Metric to U.S. Customary (Approximate)]
\begin{itemize}[left=0pt]
  \item \( 1 \text{ inch} \approx 2.54 \text{ cm} \)
  \item \( 1 \text{ meter} \approx 3.28 \text{ ft} \)
  \item \( 1 \text{ kilometer} \approx 0.62 \text{ miles} \)
  \item \( 1 \text{ kg} \approx 2.2 \text{ lb} \)
  \item \( 1 \text{ liter} \approx 1.06 \text{ qt} \)
\end{itemize}
\end{tcolorbox}




\section*{6. Complex Numbers}

\begin{tcolorbox}[title=Definition of \(i\)]
\[
i = \sqrt{-1}, \quad i^2 = -1
\]
\end{tcolorbox}

\begin{tcolorbox}[title=Powers of \(i\) Cycle]
\[
\begin{aligned}
i^1 &= i \\
i^2 &= -1 \\
i^3 &= -i \\
i^4 &= 1 \\
i^5 &= i \quad \text{(cycle repeats every 4)}
\end{aligned}
\]


\end{tcolorbox}

\newpage
\section*{8. Sequences: Arithmetic \& Geometric}

\begin{tcolorbox}[title=Arithmetic Sequence]
\[
\text{nth term: } a_n = a_1 + (n - 1)d
\]
\[
\text{Sum of } n \text{ terms: } S_n = \frac{n}{2}(a_1 + a_n)
\]
\textit{Where: } \( a_1 \) = first term, \( d \) = common difference
\end{tcolorbox}

\begin{tcolorbox}[title=Geometric Sequence]
\[
\text{nth term: } a_n = a_1 \cdot r^{n-1}
\]
\[
\text{Sum of } n \text{ terms: } S_n = a_1 \cdot \frac{1 - r^n}{1 - r} \quad (r \neq 1)
\]
\textit{Where: } \( r \) = common ratio
\end{tcolorbox}


\section*{9. Polygons: Interior and Exterior Angles}

\begin{tcolorbox}[title=Sum of Interior Angles]
\[
\text{Sum} = (n - 2) \times 180^\circ
\]
\textit{Where } \( n \) is the number of sides in the polygon.
\end{tcolorbox}

\begin{tcolorbox}[title=Each Interior Angle (Regular Polygon)]
\[
\text{Each angle} = \frac{(n - 2) \times 180^\circ}{n}
\]
\textit{Only valid if the polygon is regular (all angles equal).}
\end{tcolorbox}

\begin{tcolorbox}[title=Each Exterior Angle (Regular Polygon)]
\[
\text{Each exterior angle} = \frac{360^\circ}{n}
\]
\textit{Exterior and interior angles are supplementary: } \\
\[
\text{Interior angle} + \text{Exterior angle} = 180^\circ
\]
\end{tcolorbox}

\section*{10. GCF and LCM}

\begin{tcolorbox}[title=GCF and LCM from a Venn Diagram Example]
\[
\begin{aligned}
24 &= 2^3 \cdot 3 \\
36 &= 2^2 \cdot 3^2 \\
\\
\text{GCF}(24, 36) &= 2^2 \cdot 3 = 12 \\
\text{LCM}(24, 36) &= 2^3 \cdot 3^2 = 72
\end{aligned}
\]
\end{tcolorbox}



\section*{11. Percent Increase and Decrease}
\begin{tcolorbox}[title=Percent Increase and Decrease]
\[
\text{New Value} = \text{Original Value} \times (1 \pm \text{Rate})
\]

\textbf{Percent Increase:} use \( + \text{Rate} \) \quad
\textbf{Percent Decrease:} use \( - \text{Rate} \)

\textit{Note: Percent rate must be in decimal form (e.g., 20\% = 0.20)}
\end{tcolorbox}



\end{multicols}

\end{document}
